% !TEX program = xelatex
% ============================================================
%  مقترح بحثي — Research Proposal Template
%  نموذج مقترح بحثي مع صفحة غلاف وجدول زمني وميزانية وقائمة مراجع
%  Compile with:
%    xelatex main.tex
%    bibtex main
%    xelatex main.tex
%    xelatex main.tex
% ============================================================
%  Features:
%    - Title page with proposer and institution
%    - Summary, Background, Objectives, Methodology
%    - Timeline table and Budget table
%    - BibTeX bibliography (apalike style)
%    - Fancy headers with proposal title
% ============================================================

\documentclass[11pt,a4paper]{article}

% ---- Geometry ----
\usepackage[a4paper,margin=2.5cm,top=2.5cm,bottom=2.5cm,headheight=16pt]{geometry}

% ---- Core packages ----
\usepackage{graphicx}
\usepackage{array}
\usepackage{booktabs}
\usepackage{tabularx}
\usepackage{multirow}
\usepackage{enumitem}
\usepackage{float}
\usepackage{caption}
\usepackage{titlesec}
\usepackage{fancyhdr}
\usepackage{setspace}
\usepackage[table]{xcolor}
\usepackage{amsmath}

% ---- RTL / Arabic language support (load before hyperref) ----
\usepackage{bidi}
\usepackage[no-math]{fontspec}
\usepackage[quiet]{polyglossia}

% ---- Fonts ----
\setmainfont[Scale=1]{Amiri-Regular.ttf}
\newfontfamily\arabicfont[Script=Arabic,Scale=0.95]{Amiri-Regular.ttf}
\newfontfamily\englishfont[Scale=0.95]{TeX Gyre Termes}
\setmonofont{Amiri-Regular.ttf}
\newfontfamily\arabicfontsf{Amiri-Regular.ttf}

% ---- Languages ----
\setdefaultlanguage[calendar=gregorian,hijricorrection=1]{arabic}
\setotherlanguage[variant=british]{english}

% ---- Hyperref ----
\usepackage[breaklinks,hidelinks]{hyperref}
\hypersetup{
  pdftitle={مقترح بحثي},
  pdfauthor={اسم الباحث}
}

% ---- Captions in Arabic ----
\addto\captionsarabic{%
  \renewcommand{\contentsname}{الفهرس}%
  \renewcommand{\figurename}{الشكل}%
  \renewcommand{\tablename}{الجدول}%
  \renewcommand{\refname}{المراجع}%
}

% ---- Section formatting ----
\titleformat{\section}{\Large\bfseries}{\thesection}{0.7em}{}
\titleformat{\subsection}{\large\bfseries}{\thesubsection}{0.6em}{}
\titlespacing*{\section}{0pt}{1.4ex plus 0.4ex}{0.9ex plus 0.3ex}
\titlespacing*{\subsection}{0pt}{1.0ex plus 0.3ex}{0.7ex plus 0.2ex}

% ---- Headers / footers ----
\pagestyle{fancy}
\fancyhf{}
\fancyhead[R]{\small\itshape مقترح بحثي: عنوان مختصر}
\fancyhead[L]{\small اسم الباحث}
\fancyfoot[C]{\small\thepage}
\renewcommand{\headrulewidth}{0.3pt}
\renewcommand{\footrulewidth}{0pt}

% ---- List spacing ----
\setlist[itemize]{topsep=0.3em, itemsep=0.15em, parsep=0pt}
\setlist[enumerate]{topsep=0.3em, itemsep=0.15em, parsep=0pt}

% ---- Table column types ----
\newcolumntype{L}[1]{>{\raggedright\arraybackslash}p{#1}}
\newcolumntype{C}[1]{>{\centering\arraybackslash}p{#1}}
\newcolumntype{R}[1]{>{\raggedleft\arraybackslash}p{#1}}

% ---- Line spacing ----
\setstretch{1.4}

\begin{document}

% ============================================================
% TITLE PAGE
% ============================================================
\thispagestyle{empty}
\begin{center}
\vspace*{2cm}
{\LARGE\bfseries مقترح بحثي (Research Proposal)}\\[1.5em]
{\Large\bfseries عنوان المشروع البحثي}\\[2em]

\begin{tabular}{rl}
\textbf{مقدم المقترح (Proposer):} & اسم الباحث \\
\textbf{الجامعة (Institution):} & اسم الجامعة \\
\textbf{القسم (Department):} & اسم القسم \\
\textbf{المشرف (Supervisor):} & د. اسم المشرف \\
\textbf{المدة (Duration):} & 12 شهراً \\
\end{tabular}

\vspace{2em}
{\large \today}
\end{center}

\newpage

% ============================================================
% 1. SUMMARY
% ============================================================
\section{الملخص (Summary)}
\label{sec:summary}

% TODO: اكتب ملخص المقترح هنا (صفحة واحدة كحد أقصى).
هذا نص تجريبي للملخص. يصف المشكلة البحثية وأهميتها، والأهداف الرئيسية، والمنهجية المتبعة، والأثر العلمي المتوقع.

\section{الخلفية (Background)}
\label{sec:background}

% TODO: اكتب الخلفية العلمية ومراجعة الأدبيات هنا. استخدم \cite{bibkey}.
هذا نص تجريبي للخلفية. يستعرض الأدبيات السابقة ذات الصلة بالموضوع \cite{example_book}.

\section{الأهداف (Objectives)}
\label{sec:objectives}

% TODO: اكتب أهداف البحث هنا.
\begin{enumerate}
    \item \textbf{الهدف الأول:} وصف الهدف الأول بالتفصيل.
    \item \textbf{الهدف الثاني:} وصف الهدف الثاني بالتفصيل.
    \item \textbf{الهدف الثالث:} وصف الهدف الثالث بالتفصيل.
\end{enumerate}

\section{المنهجية (Methodology)}
\label{sec:methodology}

% TODO: اكتب منهجية البحث هنا بالتفصيل.
\subsection{مراحل البحث}
\begin{enumerate}
    \item \textbf{المرحلة الأولى:} المراجعة والأدبيات (الأشهر 1--3).
    \item \textbf{المرحلة الثانية:} جمع البيانات (الأشهر 4--7).
    \item \textbf{المرحلة الثالثة:} التحليل (الأشهر 8--10).
    \item \textbf{المرحلة الرابعة:} كتابة التقرير (الأشهر 11--12).
\end{enumerate}

% ============================================================
% TIMELINE
% ============================================================
\section{الجدول الزمني (Timeline)}
\label{sec:timeline}

% TODO: عدّل الجدول الزمني حسب مراحل مشروعك.

\begin{table}[H]
\centering
\caption{الجدول الزمني للمشروع (12 شهراً)}
\small
\definecolor{bar}{HTML}{2E6DA4}
\begin{tabular}{|l|c|c|c|c|}
\hline
\textbf{المرحلة} & \textbf{الأشهر 1--3} & \textbf{الأشهر 4--7} & \textbf{الأشهر 8--10} & \textbf{الأشهر 11--12} \\
\hline
المراجعة والأدبيات & \cellcolor{bar} & & & \\
\hline
جمع البيانات & & \cellcolor{bar} & & \\
\hline
التحليل & & & \cellcolor{bar} & \\
\hline
كتابة التقرير & & & & \cellcolor{bar} \\
\hline
\end{tabular}
\end{table}

% ============================================================
% BUDGET
% ============================================================
\section{الميزانية (Budget)}
\label{sec:budget}

% TODO: عدّل بنود الميزانية حسب احتياجات مشروعك.

\begin{table}[H]
\centering
\caption{الميزانية التفصيلية للمشروع}
\begin{tabularx}{\textwidth}{L{5cm} X R{3cm}}
\toprule
\textbf{البند} & \textbf{التفاصيل} & \textbf{التكلفة (\LR{USD})} \\
\midrule
المعدات & أجهزة وبرمجيات & \LR{2,000} \\
المواد الاستهلاكية & كواشف ومستلزمات & \LR{1,500} \\
السفر & مؤتمرات وميداني & \LR{1,000} \\
النشر & رسوم النشر المفتوح & \LR{500} \\
\midrule
\textbf{المجموع} & & \textbf{\LR{5,000}} \\
\bottomrule
\end{tabularx}
\end{table}

% ============================================================
% REFERENCES
% ============================================================
\section{المراجع (References)}
\label{sec:references}

% Add your .bib file as bibliography.bib in the project folder.
% Run: xelatex -> bibtex -> xelatex -> xelatex
\bibliography{bibliography}
\bibliographystyle{apalike}

\end{document}
